\section{Architecture \& Implementation}

The \texttt{GeometricLatticeIsomorphism} module is implemented in PyTorch, leveraging registered buffers for zero-overhead matrix multiplication.

\subsection{PyTorch Production Code}

\begin{lstlisting}[language=Python]
import math
from typing import cast
import torch
import torch.nn as nn
from nrc.math import PHI_FLOAT

class GeometricLatticeIsomorphism(nn.Module):
    """Enhancement #18: Geometric Lattice Isomorphism Projection Protocol."""

    def __init__(self, high_dim_features: int):
        super().__init__()
        self.features = high_dim_features
        self.qrt_damping_angle = math.atan(math.sqrt(PHI_FLOAT))
        self.register_buffer("isomorphism_matrix", self._build_isomorphism_matrix())

    def _build_isomorphism_matrix(self) -> torch.Tensor:
        matrix = torch.eye(self.features)
        cos_val = math.cos(self.qrt_damping_angle)
        sin_val = math.sin(self.qrt_damping_angle)

        for i in range(self.features - 1):
            if i % 2 == 0:
                matrix[i, i] = cos_val * PHI_FLOAT
                matrix[i, i + 1] = -sin_val / PHI_FLOAT
                matrix[i + 1, i] = sin_val / PHI_FLOAT
                matrix[i + 1, i + 1] = cos_val * PHI_FLOAT

        return matrix

    def forward(self, hidden_states: torch.Tensor) -> torch.Tensor:
        projected_states = torch.matmul(hidden_states, cast(torch.Tensor, self.isomorphism_matrix))
        return cast(torch.Tensor, projected_states)
\end{lstlisting}

\subsection{Execution Workflow}
\begin{enumerate}
    \item \textbf{Buffer Initialization}: `_build_isomorphism_matrix()` computes the static QRT lattice transformation matrix upon instantiation.
    \item \textbf{Registration}: `self.register_buffer` binds the matrix to model state without marking it as a trainable parameter.
    \item \textbf{Fast Matrix Projection}: `torch.matmul` projects input hidden states onto the stabilized coordinate grid during forward execution.
\end{enumerate}
